\section{Experiments}
\label{sec:experiments}

We evaluate Quake using a number of benchmarks and summarize our main findings:

\begin{enumerate}[leftmargin=*, itemsep=-1pt, topsep=0pt]
    \item Quake achieves the lowest search time across all dynamic workloads compared to state-of-the-art graph indexes, with $1.5$-$13\times$ lower search latency than HNSW, DiskANN, and SVS while having $18$-$126\times$ lower update latency.
    \item We also find that APS matches the nprobe of an oracle across recall targets on \textsc{Sift1M}, with only a $17$-$29\%$ increase in latency relative to the oracle. 
    \item APS performs on-par or better than existing early termination methods~\cite{zhang2023fast, chen_spann_nodate, li2020improving} and requires no offline tuning. %Additionally, in contrast to existing methods for early termination, APS requires no offline tuning.
    \item Quake's NUMA-aware query processing exhibits linear scalability and high memory bandwidth utilization on the \textsc{MSTuring100M} dataset. Quake achieves $20\times$ and $4\times$ lower query latency compared to single-threaded and non-NUMA aware configurations, respectively.% version and $4\times$ lower latency compared to a non-NUMA aware configuration.
\end{enumerate}

\subsection{Workloads}

We performed our evaluation on a diverse set of real-world and synthetic workloads.

\paragraph{Wikipedia-12M} 
This dataset and workload trace are derived from monthly Wikipedia page additions and page-view \cite{wiki} frequencies between April 2013 and December 2021. We consider only pages about people or those linking to people. The dataset begins with 1.6 million pages and grows to 12 million after 103 updates, and therefore the average update size is $\approx$ 100,000 vectors. Embeddings are generated by training DistMult \cite{distmult} graph embeddings (via Marius \cite{mohoney2021marius, waleffe2023mariusgnn}) on the Wikipedia link structure, and use the inner product metric.

The workload simulates monthly inserts of new pages, followed by 100,000 search queries sampling page embeddings with probability proportional to their page views corresponding to roughly a 50/50 read/write ratio. This setting imitates evolving interest and periodic growth of the dataset.

\paragraph{OpenImages-13M}
Using the methodology described by SVS~\cite{aguerrebere2024locally}, we generate a workload of 13M images from the Open Images dataset \cite{Kuznetsova_2020}. Embeddings are produced using Clip \cite{radford2021learningtransferablevisualmodels} in an inner product metric space. The workload maintains a sliding window of 2 million resident vectors and inserts and deletes vectors based on class labels until all 13 million vectors have been indexed at least once. Each insert and delete affects roughly 110K vectors. After each insert and delete operation, we run 1,000 queries randomly sampled from the entire vector set. This scenario stresses both insertion and deletion performance as well as sustained query latency.

\paragraph{Workload Generator} 
To test performance under varying workload properties, we employ a configurable workload generator applicable to any vector dataset. The key parameters are: number of vectors per operation, operation count, operation mix (read/write ratio), and spatial skew. For skewed workloads, vectors are clustered and sampled from to produce queries and updates, reflecting hot spots in the vector space.

We construct two example workloads from a 10M vector subset of the MSTuring~\cite{bigann} dataset using L2 distance:

\begin{itemize}[leftmargin=*, itemsep=-1ex, topsep=0pt]
    \item \textbf{MSTuring-RO}: A pure search workload. We uniformly sample from 100,000 provided query vectors and execute 100 search operations, each querying 10,000 vectors. This setup tests search efficiency in a static setting.
    \item \textbf{MSTuring-IH}: A dynamic workload interleaving inserts and searches. Beginning with 1 million vectors, the dataset grows to 10 million as we process 1,000 operations with a 90\% insert and 10\% search ratio. This tests the ability to handle large-scale growth while maintaining query quality.
\end{itemize}

We use the datasets \textsc{Sift1M} \cite{jegou2010product}, \textsc{Sift10M},  and \textsc{MSTuring100M} \cite{bigann} to conduct microbenchmarks.

\subsection{Experimental Setup}

\edit{Large-scale experiments are run on a 4-socket server with Intel Xeon Gold 6148 CPUs (80 cores, 160 threads), 500 GB RAM across 4 NUMA nodes, and ~300 GB/s total memory bandwidth. Some microbenchmarks (Tables \ref{tab:recall_overhead}, \ref{tab:early_term}, \ref{tab:maintenance_ablation}, and \ref{tab:two_level_results}) are run on a 2023 MacBook Pro with a M2 Max chip.} 

Search queries are processed one at a time and we report the total time to process all queries to reach a target of 90\% recall for $k=100$. \edit{Unless otherwise stated, all search numbers
use a single worker thread. Quake additionally reports a
multi-thread configuration \textbf{Quake-MT} (16 threads) where partition scans are parallelized for individual queries, while \textbf{Quake-ST} uses a single thread for search}. For updates we report the total update and maintenance time, where both Quake and the baselines process updates in batches \edit{using 16 threads}. This setup simulates an online environment where queries arrive individually, and updates are applied in batches. We report maintenance time separately from update latency, as maintenance can be conducted in the background in online systems ~\cite{xu_spfresh_2023}. 

\paragraph{Baselines:} We compare Quake against several state-of-the-art methods, including both partitioned and graph-based indexes:

\begin{itemize}[leftmargin=*, itemsep=-.25em, topsep=0pt]
    \item \textbf{Faiss-IVF}~\cite{douze2024faisslibrary}: A popular inverted file (IVF) index in Faiss. It handles updates but does no maintenance.
    \item \textbf{DeDrift}~\cite{baranchuk_dedrift_2023}: An incremental maintenance strategy designed to reduce clustering drift by periodically reclustering large partitions together with small ones. We implement DeDrift’s logic within Quake.
    \item \textbf{LIRE}~\cite{xu_spfresh_2023}: Maintenance procedure used by SpFresh. LIRE incrementally splits large clusters and deletes small clusters after updates, followed by local reassignments. We implement LIRE’s approach within Quake.
    \item \textbf{ScaNN}~\cite{guo_accelerating_2020}: A state-of-the-art highly optimized partitioned index system. It uses an unpublished incremental maintenance procedure similar to LIRE.
    \item \textbf{Faiss-HNSW}~\cite{10.1109/TPAMI.2018.2889473}: A graph-based approach (HNSW) implemented in Faiss. It supports incremental inserts but not deletes. Thus, for workloads with deletions, we omit Faiss-HNSW from those comparisons.
    \item \textbf{DiskANN}~\cite{singh2021freshdiskann}: System built around the Vamana~\cite{10.5555/3454287.3455520} index with support for dynamic updates.
    \item \textbf{SVS}~\cite{aguerrebere2024locally}: A recently released optimized implementation of the Vamana index with support for dynamic updates. 
\end{itemize}

\begin{table*}[h]
\small
\centering
\begin{tabular}{|l|cccc|cccc|ccc|cccc|}
\hline
\multirow{2}{*}{\textbf{Method}} 
& \multicolumn{4}{c|}{\textsc{Wikipedia-12M}} 
& \multicolumn{4}{c|}{\textsc{OpenImages-13M}} 
& \multicolumn{3}{c|}{\textsc{MSTuring10M-RO}} 
& \multicolumn{4}{c|}{\textsc{MSTuring10M-IH}} \\ \cline{2-16}
& \textbf{S} & \textbf{U} & \textbf{M} & \textbf{T}
& \textbf{S} & \textbf{U} & \textbf{M} & \textbf{T}
& \textbf{S} & \textbf{M} & \textbf{T}
& \textbf{S} & \textbf{U} & \textbf{M} & \textbf{T} \\ \hline
Quake-MT   & 1.53 & .01 & .44 & \textbf{1.98}  &
            .03 & .02 & .10 & \textbf{.15}   &
            .63 & .08 & .71          &
            .54 & .02 & .14 & \textbf{.70} \\ 
Quake-ST   & 9.48 & .01 & .44 & 9.93  &
            .14 & .02 & .10 & .26               &
            2.43 & .08 & 2.51                   &
            2.12 & .02 & .14 & 2.28 \\ 
Faiss-IVF  & 165.8 & .005 & 0 & $165.8^{+}$ &
            .45 & .01 & 0 & .46     &
            12.25 & 0 & 12.25          &
            13.72 & .01 & 0 & 13.73 \\ 
DeDrift    & 132.6 & .03 & .19 & $132.8^{+}$ &
            .23 & .03 & .19 & .45     &
            – & – & –                     &
            19.17 & .03 & .55 & 19.75 \\ 
LIRE       & 44.2 & .03 & .38 & $44.61^{*+}$ &
            .15 & .05 & .11 & .31     &
            – & – & –                     &
            9.08 & .02 & .21 & $9.32^{*}$ \\ 
ScaNN      & 50.27 & 1.75 & 0 & $52.02^{+}$ &
            .41 & .21 & 0 & .62     &
            2.97 & 0 & 2.97             &
            6.70 & .09 & 0 & 6.79 \\ 
Faiss-HNSW & 14.65 & .18 & 0 & 14.83 &
            – & – & – & –                 &
            1.9 & 0 & 1.9             &
            1.27 & 1.38 & 0 & 2.64 \\ 
DiskANN    & 12.11 & .32 & 0 & 12.43 &
            .22 & 1.53 & 0 & $1.75^{*}$     &
            1.16 & 0 & 1.16             &
            .81 & .48 & 0 & 1.28 \\ 
SVS        & 20.54 & .57 & 0 & $21.11^{*}$ &
            .29 & 2.32 & 0 & 2.61     &
            .33 & 0 & \textbf{.33}             &
            2.11 & .24 & 0 & $2.35^{*}$ \\ \hline
\end{tabular}
\caption{\edit{Total workload time breakdown} \textbf{in hours}.  \textbf{S}: search, \textbf{U}: update, \textbf{M}: maintenance, \textbf{T}: overall total. Recall target = 90\% and k=100. Search queries are processed one-at-time, updates are processed in batches, maintenance is conducted after each batch of search or update operations. $^{*}$Denotes the method is unable to meet the recall target with static query parameters. $^{+}$Denotes the method did not finish in a 24 hour time budget, for these we estimate the runtime based on a $10\%$ subsample of search queries.}
\label{tab:workload_comparison}
\end{table*}

% \begin{table*}[h]
% \small
% \centering
% \begin{tabular}{|l|ccc|ccc|cc|ccc|}
% \hline
% \multirow{2}{*}{\textbf{Method}} 
% & \multicolumn{3}{c|}{\textbf{Wikipedia-12M}} 
% & \multicolumn{3}{c|}{\textbf{OpenImages-13M}} 
% & \multicolumn{2}{c|}{\textbf{MSTuring10M-RO}} 
% & \multicolumn{3}{c|}{\textbf{MSTuring10M-IH}} \\ \cline{2-12} 
% & \textbf{Search} & \textbf{Update} & \textbf{Maint.} 
% & \textbf{Search} & \textbf{Update} & \textbf{Maint.} 
% & \textbf{Search} & \textbf{Maint.} 
% & \textbf{Search} & \textbf{Update} & \textbf{Maint.} \\ \hline
% \edit{Quake-MW}         & \textbf{.53 ms} & .004 ms & 7.6 s  & \textbf{1.01 ms} & .006 ms & 1.71 s    & 2.28 ms & 2.89 s & \textbf{1.94 ms} & .01 ms & .486 s \\ \hline
% \edit{Quake-SW}         & \edit{3.28 ms} & \edit{.004 ms} & \edit{7.6 s}  & - & - & -   & - & - & - & - & - \\ \hline
% Faiss-IVF      & 57.4 ms         & .002 ms & -       & 16.1 ms & .003 ms & -      & 44.1 ms          & -      & 49.4 ms          & .005 ms & - \\ \hline
% DeDrift        & 45.9 ms            & .01 ms     & 34.5 s    & 8.34 ms & .011 ms & 3.43 s    & -                & -      & 69.0 ms          & .01 ms & 1.98 s \\ \hline
% LIRE           & 15.3* ms            & .01 ms     & 6.64 s    & 5.34 ms & .016 ms & 1.90 s    & -                & -      & 32.7* ms          & .009 ms & 0.77 s \\ \hline
% ScaNN          & 17.4 ms            & .60 ms     & -       & 14.8 ms & .07 ms & -      & 10.7 ms             & -      & 24.1 ms             & .037 ms    & - \\ \hline
% Faiss-HNSW     & 5.07 ms         & .06 ms  & -       & -    & -    & -      & 6.83 ms             & -      & 4.57 ms             & .55 ms    & - \\ \hline
% DiskANN        & 4.19 ms         & .11 ms  & -       & 7.99* ms & .5 ms & -    & 4.17 ms          & -      & 2.9 ms           & .19 ms & - \\ \hline
% SVS            & 7.11* ms            & .197 ms     & -       & 10.34 ms &  .758ms & -      & \textbf{1.18 ms}             & -      & 7.61* ms             & .095 ms    & - \\ \hline
% \end{tabular}
% \caption{Workload performance: showing \emph{single-query} search latency to reach a recall of .9, mean single-update latency, and mean maintenance latency per operation. \edit{All baselines and Quake-SW process queries with a single thread, Quake-MW uses 16 threads for intra-query parallelism}  *Denotes the config is unable to meet the recall target with static query parameters.} 

% \label{tab:mean_latency_combined}
% \end{table*}

We configure the main parameters of Quake and the baselines as follows. We disable vector quantization/compression for all baselines, as not all baselines support it. For partitioned indexes we use $sqrt(|X_0|)$ partitions where $|X_0|$ is the initial number of vectors in the workload. For the graph indexes, we use a graph degree of 64. For LIRE and Quake, we set the partition refinement radius $r=50$. For Quake we use a single level of partition, set $\tau=250ns$, use one iteration of k-means for refinement and set $f_M$ between 1\%-10\%. All systems use 16 threads for updates and maintenance (if applicable). SCANN, DiskANN, and SVS perform maintenance eagerly during an update, therefore we do not measure maintenance time separately from update time. We consider maintenance after each operation for all methods. Throughout all experiments, indexes search parameters are tuned to achieve an average of 90\% recall for $k=100$ across the workloads. 

\subsection{End-to-End Evaluation}

\paragraph{Comparison with Baselines} Table~\ref{tab:workload_comparison} shows that Quake consistently achieves lower search, update and total time on all workloads. On the \textsc{Wikipedia-12M} workload, where the dataset grows over time and partitions can become unbalanced, the multi-threaded \edit{Quake-MT} takes 1.53 hours to process searches, \edit{while single-threaded Quake-ST takes 9.48 hours}. In contrast, Faiss-IVF climbs to 165 hours due to the lack of maintenance, DeDrift reaches 132 hours despite its rebalancing efforts, LIRE is unable to meet the recall target and takes 44 hours and SCANN performs similarly with poor update latency due to over-eager maintenance applied during updates. Even the best-performing graph-based method, DiskANN, takes 12 hours. Thus, \edit{Quake-MT} is $8\times$ faster to search than the strongest baseline on this workload derived from real-world access patterns.

On the \textsc{OpenImages-13M} workload, which includes both insertions and deletions, \edit{Quake's} multi-threaded and single-threaded search times are .03 and .14 hours respectively. The best competing approach, DiskANN, records .22 hours, making Quake-MT $7.3\times$ and Quake-ST $1.6\times$ faster. Faiss-HNSW does not support deletions so it is omitted. Both SVS's and DiskANN’s delete consolidation is expensive, leading to orders of magnitude higher update latency than partitioned indexes, illustrating that graph-based indexes struggle with dynamic operations. Quake’s continuous maintenance keep partitions balanced, achieving low latency and stable recall.

For the static, read-only, \textsc{MSTuring10M-RO} workload, Quake’s maintenance improves the index structure even without data changes, adapting partitions to the query pattern. For \edit{Quake-MT}, this yields a search time of .63 hours and \edit{Quake-ST} takes 2.43 hours to conduct the search. However, the \textsc{MSTuring10m} dataset is especially challenging for partitioned indexes, as they need to scan roughly $10\%$ of all partitions in order to meet the recall target. In contrast, the well-optimized SVS library exhibits a superior search time of .33 hours, demonstrating that in static settings, well-optimized graph indexes are strong competition.

On \textsc{MSTuring10M-IH}, where the dataset grows from one to ten million vectors, \edit{Quake-MT} achieves a total search time of .54 hours. DiskANN, the second-best performer, has a search time of .81 hours, making Quake-MT $1.5\times$ faster due to intra-query parallelism. However, single-threaded Quake is $2.6\times$ slower than DiskANN, further illustrating the search efficiency of graph indexes. The other baselines fail to maintain the recall target or suffer from high latency due to their static parameters and inability to prevent partition skew.

Overall, these results demonstrate that Quake’s combination of adaptive partition scanning, incremental maintenance, and NUMA-aware parallelism consistently delivers low-latency queries at the desired recall. Systems without maintenance (Faiss-IVF) suffer from skew-induced latency increases, those tied to static search parameters (LIRE) struggle to maintain recall without incurring higher query times, and graph-based methods (Faiss-HNSW, DiskANN) face substantial overheads when handling updates and deletions. By integrating these components, Quake matches the low update cost of partitioned indexes while outperforming graph indexes in search latency in dynamic workloads. Our design is an advancement to the state-of-the-art, providing stable, efficient performance across diverse, and evolving workloads.

\begin{figure}[h]
\centering
\includegraphics[width=0.43\textwidth]{sections/figures/stacked_latency_recall_partitions.pdf}
\caption{Comparison of single-threaded search latency, recall and number of partitions for Quake vs. maintenance approaches LIRE and DeDrift on \textsc{Wikipedia-12M}. Quake maintains stable latency and recall throughout the workload.}
\label{fig:partitioned_maintenance}
\end{figure}


\paragraph{Comparison with Partitioned Index Maintenance Methods} 
Here we perform a detailed comparison with LIRE and DeDrift, measuring the latency, recall, and number of partitions over time on the \textsc{Wikipedia-12M} workload. For a fair comparison, we use a single-thread to highlight the advantages of APS and maintenance in Quake. The results are shown in Figure \ref{fig:partitioned_maintenance}. First looking at recall, we see that Quake maintains a stable recall of near 90\%, while LIRE's recall degrades over time as it uses a static nprobe. DeDrift's recall stays relatively constant, as it does not adjust the number of partitions and therefore does not need to adjust nprobe. However, when turning our attention to latency, we see that Quake has near-constant stable latency, even as the dataset grows, while DeDrift's latency increases significantly with time. In terms of the number of partitions, we see DeDrift stays constant while Quake and Lire increase by $2.5\times$ and $10\times$ respectively. LIRE uses significantly more partitions because it uses size thresholding to determine when to split, regardless of whether a given partition is hot or not. Quake on the other hand only splits partitions if their contribution to the cost model is high, allowing for more efficient maintenance. These results show that Quake's approach to maintenance is superior to existing methods for partitioned index maintenance in minimizing query latency and recall stability.

\begin{table}[h]
\footnotesize
\centering
\caption{Ablation Study on \textsc{Wikipedia-12M} showing mean search latency and the standard deviation of recall.}
\label{tab:ablation_study}
\begin{tabular}{|l|c|c|}
\hline
\textbf{Configuration} & \textbf{Search Latency} & \textbf{Recall Std.} \\ \hline
Quake-MT            & 0.53 ms      &  .008              \\ \hline
Quake-MT w/o APS            & 0.50 ms     &  .025                \\ \hline
Quake-ST           & 3.28 ms    &  .005                 \\ \hline
Quake-ST w/o APS       & 3.18 ms      &  .025             \\ \hline
Quake-ST w/o Maint/APS & 45.20 ms   &  .014                   \\ \hline
\end{tabular}
\end{table}

\paragraph{Wikipedia-12M Ablation}
To quantify the contributions of Quake components, we disabled key features and measured the impact on \textsc{Wikipedia-12M} workload in Table \ref{tab:ablation_study}.  We see that disabling APS has little impact on the query latency, as Quake can achieve a low latency even in the static nprobe setting. However, APS provides significantly more recall stability, as evidenced by the increase in standard deviation when APS is disabled. Disabling NUMA-aware multi-threading, however, shows a $6\times$ increase in query latency, demonstrating the benefit of parallelization of partition scans. Finally, we disable maintenance and see a significant increase in latency, similar to the latency of Faiss-IVF; here partitions become extremely imbalanced due to the skew in the workload (see Figure \ref{fig:wiki_combined}) causing queries to scan more vectors and therefore increasing latency. This further demonstrates the necessity for maintenance for dynamic workloads. In conclusion, each piece of Quake contributes to its performance in terms of both recall stability and minimal query latency.

\begin{figure}[h]
\centering
\includegraphics[width=0.4\textwidth]{sections/figures/qps_vs_batchsize.pdf}
\caption{Multi-query evaluation on \textsc{Wikipedia-12M} with 10,000 search queries. QPS @ recall=90\% is measured for all baselines while varying the batch size. All methods use 16 threads to process queries.}
\label{fig:qps_vs_batch_size}
\end{figure}


\edit{\subsection{Multi-Query Processing}}
Here we compare the search performance of Quake and baseline methods in a static batched query setting. Figure \ref{fig:qps_vs_batch_size} shows the QPS at a recall of 90\%, varying the number of queries in a batch on the \textsc{Wikipedia-12M} workload. The dataset includes all 12M vectors, with 10,000 queries sampled according to Wikipedia page views from December 2021. For Quake, FaissIVF, and SCANN, we use 3,500 partitions; for FaissHNSW, SVS, and DiskANN, we set the graph degree to 64. All experiments use 16 threads for query processing. Quake employs the multi-query execution policy in \cite{mohoney_high-throughput_2023} and \cite{pound2025micronn}, grouping queries by the partitions they access and scanning each partition exactly once per batch in parallel.

Quake consistently outperforms all baselines across every batch size, with an increasing advantage as batch sizes grow. At the largest batch size (10,000 queries), Quake achieves a $6.7\times$ speedup over FaissIVF and SCANN. This performance gain stems from Quake's efficient multi-query execution strategy, where it scans each partition once per batch, in contrast to FaissIVF and SCANN, which scan partitions individually per query. Compared to DiskANN, the strongest graph-based competitor, Quake still maintains a substantial $1.8\times$ speedup. These results demonstrate that Quake delivers high performance not only in single-query scenarios but also in large multi-query workloads.

%  showing that Quake exhibits superior performance even with limited resources.
% explosion
% exceptional 

\subsection{Scalability}

We tested Quake’s parallel scalability by varying the number of threads. In Figure~\ref{fig:scalability_threads} we measure the mean search latency and scan throughput (bytes scanned / query latency) on \textsc{MSTuring100M} to reach a recall of 90\%. Note that this dataset has 100 million vectors and is $10\times$ larger than the datasets we compared against previously. We compare our NUMA-aware parallelism with one in which NUMA is disabled. For both configurations, we see near linear scalability up to around 8 workers, where the non-NUMA latency performs best (28ms). The NUMA configuration however further improves and at 64 workers achieves a latency of 6ms. Looking at the scan throughput, we see that NUMA achieves a peak throughput of 200GBps. We do not completely saturate memory bandwidth due to other overheads involved in query processing (topk sorting, memory allocations, coordination). In conclusion, NUMA-aware intra-query parallelism is an effective mechanism for decreasing query latency by utilizing the full memory capabilities of multi-core machines.

\begin{figure}[ht!]
    \centering
    \begin{subfigure}{0.23\textwidth}
        \centering
        \includegraphics[width=\textwidth]{sections/figures/latency_vs_num_workers.pdf}
        \caption{Mean Search Latency}
    \end{subfigure}
    \hfill
    \begin{subfigure}{0.23\textwidth}
        \centering
        \includegraphics[width=\textwidth]{sections/figures/throughput_vs_num_workers.pdf}
        \caption{Scan Throughput}
    \end{subfigure}
    \caption{\textsc{MSTuring100M}: Scaling the number of threads with and without NUMA.}
    \label{fig:scalability_threads}
\end{figure}

% \paragraph{Maintenance Frequency:} 
% We evaluated the effect of different maintenance frequencies. Increasing maintenance frequency slightly raises update costs but significantly lowers query latency. 

% \paragraph{Effect of Insert-Delete Ratio:} 
% Table~\ref{tab:insert_delete_ratio} examines performance under varying insert/delete ratios on MSTuring10M. Quake handles these changes gracefully, maintaining low query latency and efficient updates even as the workload composition shifts.

% \begin{table}[h]
% \centering
% \caption{Impact of Insert-Delete Ratio on MSTuring10M. Quake maintains stable performance even as the update mix changes.}
% \label{tab:insert_delete_ratio}
% \begin{tabular}{|l|c|c|}
% \hline
% \textbf{Insert-Delete Ratio} & \textbf{Query Latency (ms)} & \textbf{Update Latency (ms)} \\ \hline
% 100\% Inserts / 0\% Deletes  & 2.01 ms                     & 0.1 ms                        \\ \hline
% 90\% Inserts / 10\% Deletes  & X ms                        & U ms                          \\ \hline
% 50\% Inserts / 50\% Deletes  & X ms                        & U ms                          \\ \hline
% 10\% Inserts / 90\% Deletes  & X ms                        & U ms                          \\ \hline
% \end{tabular}
% \end{table}

% \subsection{Varying Recall Targets with APS}
% We compared APS to an oracle that knows the exact number of partitions (nprobe) required for each query. As shown in Figure~\ref{fig:adaptive_nprobe_vs_oracle}, APS closely matches the oracle’s nprobe, recall, and latency, across all recall targets, proving APS is effective despite its lack of prior knowledge. Because of the precomputation and recomputation performance optimizations (Table \ref{tab:recall_overhead}), the recall estimation in APS introduces minimal overheads which is evidenced by the query latency, which stays within 20\% of the oracle. 

\begin{table}[h]
\footnotesize
\centering
\caption{\edit{Early-termination methods on \textsc{Sift1M} with a partitioned index with 1000 partitions. Each row shows the average recall, nprobe, and 
mean per-query latency in milliseconds over 10000 queries after tuning for a specific recall target for $k=100$. We also report the total offline tuning time in seconds, where APS needs \emph{no} offline
tuning.}}
\label{tab:early_term}
\begin{tabular}{lcccccc}
\toprule
Method & Target & Recall & $nprobe$ & Latency & Offline Tuning \\
\midrule
APS         & 80\% & 82.1\% & 11.8 & 0.34ms & \textbf{0} \\
            & 90\% & 91.2\% & 20.2 & 0.48ms & \textbf{0} \\
            & 99\% & 98.9\% & 50.1 & 0.96ms & \textbf{0} \\
\addlinespace
Auncel \cite{zhang2023fast}      & 80\% & 85.7\% & 16.4 & 0.41ms & 66.3s \\
            & 90\% & 98.1\% & 73.8 & 1.29ms & 73.8s \\
            & 99\% & 99.7\% & 95.9 & 1.61ms & 83.2s \\
\addlinespace
SPANN \cite{chen_spann_nodate}      & 80\% & 81.6\% & 11  & 0.31ms & 173s \\ 
            & 90\% & 90.2\% & 19  & 0.43ms & 183s \\
            & 99\% & 99.0\% & 70  & 1.07ms & 259s \\
\addlinespace
LAET \cite{li2020improving}       & 80\% & 81.3\% & 10.5 & 0.29ms & 81s \\
            & 90\% & 90.5\% & 18.2 & 0.42ms & 104s \\
            & 99\% & 99.0\% & 58.3 & 1.03ms & 232s \\
\addlinespace
Fixed & 80\% & 81.7\% & 11  & 0.33ms & 318s \\
            & 90\% & 90.3\% & 19  & 0.44ms & 330s \\
            & 99\% & 99.0\% & 65  & 1.16ms & 424s \\
\addlinespace
Oracle      & 80\% & 83.3\% & 11.5 & 0.29ms & 320s \\
            & 90\% & 92.4\% & 19.3 & 0.41ms & 331s \\
            & 99\% & 99.2\% & 42.0 & 0.74ms & 368s \\
\bottomrule
\end{tabular}
\end{table}

\edit{\subsection{Comparison with Early Termination Methods}}
Table \ref{tab:early_term} compares early-termination methods on \textsc{Sift1M}, highlighting the tradeoff between query latency, tuning time, and recall, we do not include ground truth generation time in the tuning time for the baselines. APS analytically estimates recall at query time, \textbf{eliminating offline tuning entirely}, while achieving latency within 30\% of the oracle across all recall targets. Fixed selects a static $nprobe$ per target via an expensive offline binary search (up to 424s), and SPANN similarly performs a binary search  and tunes a centroid-distance threshold; both closely match recall targets but incur higher latency at 90\% and 99\% recall. LAET trains a per-query prediction model, incurring moderate tuning overhead (81–232s), and matches recall targets with slightly higher latency compared to APS. Auncel is the most similar method to APS, as it aims to analytically estimate recall using partition intersection volumes, however its volume estimation requires calibration, and it is a conservative method, overshooting recall. We tune Auncel by binary searching a geometric parameter ($a$), overshooting recall significantly (up to 8.1 pp) and increasing latency by up to 169\% compared to APS. Finally, the Oracle, which scans the minimal amount of partitions per-query, serves as a practical lower bound on achievable latency, though with prohibitively high tuning cost. The tuning overhead of the baselines demonstrates a significant burden for online scenarios where the queries, data, and index change. The tuning burden worsens at scale, where ground truth generation cost grows linearly with the data size, and running queries multiple times to conduct binary search takes longer. APS thus provides near-optimal performance without tuning overhead, matching or exceeding all baselines.

\begin{table}[ht]
\footnotesize
\centering
\caption{
\textsc{Sift10M}: Recall and per‐level search latency ($\ell_{0}$, $\ell_{1}$, total) for a single‐level baseline (L$_0$: 40,000 partitions; L$_1$: 1 partition) and two‐level index (L$_0$: 40,000 partitions; L$_1$: 500 partitions), where recall targets are varied at each level. The single‐level baseline is the first row of each $\tau_{r}(0)$ block.
}
\label{tab:two_level_results}
\begin{tabular}{cccccc}
\toprule
$\tau_{r}(0)$ & $\tau_{r}(1)$ & Recall & $\ell_{0}$ (ms) & $\ell_{1}$ (ms) & Total (ms) \\
\midrule
\multirow{7}{*}{\centering 80\%} 
  & \textemdash & 81.2\% & 2.07 & 4.85 & 6.92 \\
\cmidrule(lr){2-6}
  & 80\%       & 74.8\% & 1.44 & 0.72       & 2.16 \\
  & 90\%       & 78.3\% & 1.56 & 1.19       & 2.75 \\
  & 95\%       & 80.1\% & 1.67 & 1.69       & 3.37 \\
  & 99\%       & 81.0\% & 1.75 & 2.57       & 4.33 \\
  & 100\%      & 81.1\% & 1.82 & 3.81       & 5.63 \\
\midrule
\multirow{7}{*}{\centering 90\%} 
  & \textemdash & 91.3\% & 2.85 & 4.89 & 7.86 \\
\cmidrule(lr){2-6}
  & 80\%       & 84.1\% & 2.07 & 0.77       & 2.84 \\
  & 90\%       & 88.2\% & 2.26 & 1.24       & 3.50 \\
  & 95\%       & 90.1\% & 2.38 & 1.72       & 4.10 \\
  & 99\%       & 91.0\% & 2.48 & 2.60       & 5.08 \\
  & 100\%      & 91.2\% & 2.62 & 3.88       & 6.50 \\
\midrule
\multirow{7}{*}{\centering 99\%} 
  & \textemdash & 99.0\% & 4.82 & 5.5 & 10.3 \\
\cmidrule(lr){2-6}
  & 80\%       & 91.4\% & 4.11 & 0.77       & 4.88 \\
  & 90\%       & 96.0\% & 4.59 & 1.27       & 5.86 \\
  & 95\%       & 97.7\% & 4.80 & 1.75       & 6.55 \\
  & 99\%       & 98.7\% & 5.08 & 2.65       & 7.74 \\
  & 100\%      & 98.9\% & 5.28 & 3.93       & 9.21 \\
\bottomrule
\end{tabular}
\end{table}


\subsection{Multi-Level Recall Estimation}
Here we evaluate the effectiveness of APS in a two-level partitioned index by measuring the impact of varying per-level recall targets on overall recall. Table~\ref{tab:two_level_results} reports results on \textsc{Sift10M} using 40{,}000 partitions at level $L_0$ and 500 partitions at level $L_1$. We use an initial search fraction of $f_M = 1.5\%$ at $L_0$ and 25\% at $L_1$.

We observe that overly aggressive early termination at $L_1$ by setting $\tau_r(1)$ too low leads to a clear degradation in total recall. For instance, at $\tau_r(0)=90\%$, reducing $\tau_r(1)$ from 99\% to 80\% lowers overall recall from 91.0\% to 84.1\%. This confirms that accurate recall estimation at the upper level is necessary to maintain accurate end-to-end recall estimates. These results justify our design decision in Section~\ref{sec:aps-algo} to fix $\tau_r(1)=99\%$ in multi-level configurations so that only the base recall target $\tau_r(0)$ needs to be provided.

In addition, the two-level index substantially reduces centroid scanning overhead. The single-level baseline must evaluate distances to all 40{,}000 centroids per query. In contrast, the two-level configuration performs an approximate search over the centroids. For example, at a recall target of $\tau_r(0) = 90\%$ with $\tau_r(1) = 99\%$ the total query latency drops from 7.86ms to 5.08ms, a 35\% reduction, driven mainly by the drop in $L_1$ latency from 4.89ms to 2.60ms.

\edit{\subsection{Maintenance Ablation}}
\label{sec:abl_maintenance}

To understand the effectiveness of the primary components of adaptive incremental maintenance (cost-model, partition refinement, and rejection), we replay a dynamic \textsc{Sift1M} trace (30\% inserts, 20\% deletes, 50\% queries) with different components disabled. All methods use a single-thread and search using APS with $k{=}100$ and a 90\% recall target. We also include LIRE as a baseline. Table \ref{tab:abl_sift} reports cumulative times in seconds. For configurations with refinement we use a refinement radius of $r_f=50$. The full Quake policy delivers the lowest search cost (86s) while meeting the recall target. If we keep the cost model but skip refinement (NoRef), maintenance time decreases significantly from 21s to 5s, yet recall slips by 2.4 pp and the search time increase by 15.4s. This shows that while refinement is the dominant cost in maintenance, it is necessary for minimizing search latency. Disabling the cost model and instead using size-based thresholding (NoCost) shows why naive size thresholds are inadequate: search time rises 8\% despite similar maintenance effort. The rejection mechanism is critical; once removed (NoRej), recall collapses to 66\% even though search and maintenance appear cheap. LIRE, which relies solely on size thresholding, is 17\% slower in search latency, confirming that the cost model, rejection mechanism, and partition refinement are essential for maintaining both index performance and quality.

\begin{table}[h]
\footnotesize
\centering
\caption{\edit{Maintenance ablation on the \textsc{Sift1M} workload. Times are cumulative (in seconds) over the course of the workload. Recall is averaged over all queries.}}
\label{tab:abl_sift}
\begin{tabular}{lrrrr}
\toprule
Maintenance Variant & Search & Update & Maint. & Recall \\
\midrule
Quake (Full)           & 86.3s & 21.7s & 21.4s & 90.5\% \\
NoRef            & 101.7s & 22.2s &  5.2s & 88.1\% \\
NoRef+NoRej      & 85.5s  & 21.0s &  1.0s & 73.0\% \\
NoRej            & 84.2s  & 19.6s & 18.5s & 66.2\% \\
NoCost           & 93.5s  & 20.0s & 20.4s & 90.1\% \\
NoCost+NoRef     & 100.7s & 21.2s &  0.8s & 87.9\% \\
\midrule
LIRE             & 100.5s & 21.2s & 11.9s & 90.0\% \\
\bottomrule
\label{tab:maintenance_ablation}
\end{tabular}
\end{table}

% \begin{figure}[h]
% \centering
% \includegraphics[width=0.4\textwidth]{sections/figures/mean_line_plots_stacked.pdf}
% \caption{APS vs. Oracle on SIFT1M. APS achieves similar recall and latency w.r.t. an oracle tuned for each recall target.}
% \label{fig:adaptive_nprobe_vs_oracle}
% \end{figure}

% \paragraph{Maintenance Frequency}

% \paragraph{Recall and Latency Percentiles:} 
% Table~\ref{tab:percentiles} reports latency percentiles on the Wikipedia-12M workload. Quake exhibits low variance, ensuring predictable performance. Even at higher percentiles, query latency remains close to the mean, reflecting robust and stable performance under dynamic conditions.

% \begin{table}[h]
% \small
% \centering
% \caption{Latency (ms) and Recall (\%) Percentiles on Wikipedia-12M}
% \label{tab:percentiles}
% \begin{tabular}{|l|ccc|}
% \hline
% \textbf{Percentile} & \textbf{Quake} & \textbf{SCANN} & \textbf{DiskANN} \\ \hline
% \multicolumn{4}{|c|}{\textbf{Latency (ms)}} \\ \hline
% 50\% (Mean)         & 2.6            & Y             & Z               \\ 
% 90\%                & X              & Y             & Z               \\ 
% 95\%                & X              & Y             & Z               \\ 
% 99\%                & X              & Y             & Z               \\ \hline
% \multicolumn{4}{|c|}{\textbf{Recall (\%)}} \\ \hline
% 50\% (Mean)         & A              & B             & C               \\ 
% 90\%                & A              & B             & C               \\ 
% 95\%                & A              & B             & C               \\ 
% 99\%                & A              & B             & C               \\ \hline
% \end{tabular}
% \end{table}
